\documentclass[11pt]{article}
\usepackage{amsmath, amsthm, amssymb}
\usepackage[margin=1.1in]{geometry}
\usepackage{hyperref}

\newtheorem{theorem}{Theorem}
\newtheorem{proposition}[theorem]{Proposition}
\newtheorem{corollary}[theorem]{Corollary}
\newtheorem{lemma}[theorem]{Lemma}
\theoremstyle{definition}
\newtheorem*{conjecture}{Conjecture}
\theoremstyle{remark}
\newtheorem*{remark}{Remark}

\newcommand{\Frob}{\operatorname{Frob}}
\newcommand{\Ind}{\operatorname{Ind}}
\newcommand{\Res}{\operatorname{Res}}
\newcommand{\Hom}{\operatorname{Hom}}
\newcommand{\sgn}{\operatorname{sgn}}
\newcommand{\ch}{\operatorname{ch}}
\newcommand{\triv}{\mathrm{triv}}
\newcommand{\fdeg}{\widetilde{f}}
\newcommand{\qbinom}[2]{\left[{#1 \atop #2}\right]_{\!t}}
\newcommand{\CC}{\mathbb{C}}
\newcommand{\ZZ}{\mathbb{Z}}

\title{$[2r-1]_t$-divisibility of $S_r$-isotypic multiplicities of\\$R_{(k^r)}$ for $2 \le k \le 2r-2$:\\a Molien-Springer proof}
\author{Clio}
\date{2026-07-31}

\begin{document}
\maketitle

\begin{abstract}
Fix integers $r \ge 2$, $k \ge 2$ and set $n = rk$, $d = 2r - 1$,
$\alpha = (k^r)$. Let $R_\alpha = R_n^{S_\alpha}$ be the parabolic
coinvariant algebra of $S_n$ (a graded $S_r$-module via the residual
block-permutation action), and let
$m_\pi(t) = \sum_\lambda \fdeg_\lambda(t) \langle s_\pi[h_k], s_\lambda\rangle$
be the graded $S_r$-isotypic multiplicities (Theorem~D).  We prove:

\begin{theorem*}
Assume $d = 2r - 1$ is prime and $2 \le k \le d - 1$. Then
\[
[d]_t \mid m_\pi(t) \qquad \text{for every } \pi \vdash r.
\]
\end{theorem*}

The proof uses the Molien formula for $R_n^{S_\alpha}$-traces of the
block-permuting elements, combined with the elementary observation that
in the wreath product $S_k \wr S_r$ the cycle lengths of any element
lie in $\bigcup_{\ell \le r} \ell \cdot \{1, 2, \ldots, k\}$; when
$k < d$ and $d$ is prime, no such length is divisible by $d$, and the
Springer-Molien numerator kills every term of the trace sum.

The theorem covers all empirically verified positive cases with
$k$ in the first period (through $r = 6$, per
\texttt{\~{}/projects/probes/2026-07-31-2r-1-divisibility/results.md}),
and proves the entire ``$k \bmod d \in \{2, \ldots, d-1\}$'' portion
of the conjecture (\cite{PROVE-2026-07-31}) restricted to $k < d$.
\end{abstract}

\section{Setup and reduction to characters}

The action of $S_r$ on $R_\alpha$ is by graded ring automorphisms, so
$R_\alpha$ splits into $S_r$-isotypic components
$R_\alpha = \bigoplus_{\pi \vdash r} V^\pi \otimes M_\pi$, with each
$M_\pi$ a graded vector space and $m_\pi(t) := \sum_i t^i \dim (M_\pi)_i$.
The graded character trace of $\sigma \in S_r$ is
\[
P_\sigma(t) \;:=\; \sum_{\pi \vdash r} \chi^\pi(\sigma) \, m_\pi(t),
\]
so by orthogonality
\[
m_\pi(t) \;=\; \frac{1}{r!} \sum_{\sigma \in S_r} \chi^\pi(\sigma) P_\sigma(t).
\]

Consequently:
\begin{lemma}\label{lem:reduce}
For any $\zeta \in \CC$, $m_\pi(\zeta) = 0$ for every $\pi \vdash r$ if
and only if $P_\sigma(\zeta) = 0$ for every $\sigma \in S_r$.
\end{lemma}

\begin{proof}
The character table of $S_r$ is invertible, so the linear map
$(m_\pi(\zeta))_\pi \leftrightarrow (P_\sigma(\zeta))_\sigma$ is a bijection.
\end{proof}

Since $d = 2r - 1$ is prime by assumption, $[d]_t = \Phi_d(t)$ is
irreducible over $\ZZ[t]$, and $[d]_t \mid m_\pi(t)$ iff $m_\pi(\zeta_d) = 0$
for $\zeta_d = e^{2\pi i/d}$. So the theorem reduces to:

\begin{proposition}\label{prop:main}
For $r \ge 2$, $d = 2r - 1$ prime, and $2 \le k \le d - 1$:
$P_\sigma(\zeta_d) = 0$ for every $\sigma \in S_r$.
\end{proposition}

\section{Molien formula for $P_\sigma(t)$}

Fix $\sigma \in S_r$ and let $\tilde\sigma \in S_n$ denote its
block-permutation lift (viewing $S_r \hookrightarrow N_{S_n}(S_\alpha) / S_\alpha$).
By Chevalley--Shephard--Todd applied to the coinvariant algebra
$R_n = \CC[V]/\langle e_1, \ldots, e_n\rangle$ of $S_n$ on $V = \CC^n$, the
graded trace of any $g \in S_n$ on $R_n$ is
\[
\text{tr}(g \mid R_n; t) \;=\; \frac{\prod_{i=1}^{n}(1 - t^i)}{\det(1 - tg \mid V)}.
\]
For $g \in S_n$ with cycle lengths $c_1, c_2, \ldots$,
$\det(1 - tg|_V) = \prod_j (1 - t^{c_j})$.  Averaging over $S_\alpha$
gives the trace on the $S_\alpha$-invariant subalgebra:

\begin{lemma}\label{lem:molien}
For every $\sigma \in S_r$ with block-lift $\tilde\sigma \in S_n$,
\[
P_\sigma(t) \;=\;
\frac{1}{|S_\alpha|} \sum_{h \in S_\alpha}
\frac{\prod_{i=1}^{n}(1 - t^i)}{\prod_j (1 - t^{c_j(\tilde\sigma h)})},
\]
where the inner product ranges over cycles of $\tilde\sigma h$ as an
element of $S_n$.
\end{lemma}

\begin{proof}
Standard: since $\tilde\sigma$ normalises $S_\alpha$,
$P_\sigma(t) = \text{tr}(\tilde\sigma \mid R_n^{S_\alpha}; t)
= |S_\alpha|^{-1} \sum_h \text{tr}(\tilde\sigma h \mid R_n; t)$;
apply the CST trace formula to each summand.
\end{proof}

\section{Order of vanishing at $\zeta_d$}

Given $g \in S_n$, let
\[
m_d(g) \;:=\; \#\{j : d \mid c_j(g)\}.
\]

\begin{lemma}\label{lem:order}
For any $g \in S_n$, the individual trace
$\text{tr}(g \mid R_n; t)$ has a zero of order exactly
$\lfloor n/d \rfloor - m_d(g)$ at $t = \zeta_d$.  In particular,
\[
\text{tr}(g \mid R_n; \zeta_d) \;=\; 0
\quad \iff \quad m_d(g) < \lfloor n/d \rfloor.
\]
\end{lemma}

\begin{proof}
The numerator $\prod_{i=1}^{n}(1 - t^i)$ vanishes at $\zeta_d$ to order
$\lfloor n/d \rfloor$ (once for each $i \in \{d, 2d, \ldots, d\lfloor n/d\rfloor\}$).
The denominator $\prod_j (1 - t^{c_j})$ vanishes to order $m_d(g)$.
The ratio is a polynomial in $t$, so the total order at $\zeta_d$ is
$\lfloor n/d\rfloor - m_d(g) \ge 0$.
\end{proof}

\section{Wreath cycle bound and proof of Proposition~\ref{prop:main}}

\begin{lemma}[Wreath cycle formula]\label{lem:wreath-cycles}
Let $\sigma \in S_r$ and $h = (h_1, \ldots, h_r) \in S_\alpha = S_k^r$,
and let $\tilde\sigma h \in S_n$ be the resulting product.  For each
cycle $(b_1\ b_2\ \cdots\ b_\ell)$ of $\sigma$, form the ``companion''
permutation $g_c := h_{b_\ell} h_{b_{\ell-1}} \cdots h_{b_1} \in S_k$.
Then the cycles of $\tilde\sigma h$ on the $\ell k$ elements of blocks
$b_1, \ldots, b_\ell$ have lengths precisely
\[
\{\ell m : m \text{ is a cycle length of } g_c\}.
\]
\end{lemma}

\begin{proof}
Classical; see e.g., Macdonald, \emph{Symmetric Functions and Hall
Polynomials} (2nd ed.), Ch.~I App.~B, or James--Kerber
\emph{Representation Theory of the Symmetric Group} Theorem 4.2.8. In
brief: trace an element $(b_1, x) \in b_1$ under repeated application of
$\tilde\sigma h$. Since $h$ acts within blocks and $\tilde\sigma$
permutes blocks, after $\ell$ applications we return to $b_1$ with
$x$ mapped by $g_c$; the return time to $x$ itself is then
$\ell \cdot (\text{order of } x \text{ under } g_c)$.
\end{proof}

\begin{corollary}\label{cor:no-d-cycle}
If $d$ is prime with $d > r$ and $k < d$, then $m_d(\tilde\sigma h) = 0$
for every $\sigma \in S_r$ and every $h \in S_\alpha$.
\end{corollary}

\begin{proof}
By Lemma~\ref{lem:wreath-cycles}, every cycle length of $\tilde\sigma h$
has the form $\ell m$ with $\ell \le r$ a cycle length of $\sigma$ and
$1 \le m \le k$ a cycle length of a companion permutation in $S_k$. Since
$d$ is prime and $\ell \le r < d$, $\gcd(d, \ell) = 1$, so
$d \mid \ell m \iff d \mid m$. But $m \le k < d$ forbids $d \mid m$.
\end{proof}

\begin{proof}[Proof of Proposition \ref{prop:main}]
Assume $2 \le k \le d - 1$. Then $n = rk \ge 2r > d$ (since $r \ge 2$ and
$d = 2r - 1$), so $\lfloor n/d \rfloor \ge 1$. By
Corollary~\ref{cor:no-d-cycle}, $m_d(\tilde\sigma h) = 0 < \lfloor n/d \rfloor$
for every $\sigma \in S_r$ and every $h \in S_\alpha$.

By Lemma~\ref{lem:order}, every summand of
Lemma~\ref{lem:molien} vanishes at $t = \zeta_d$. Hence
$P_\sigma(\zeta_d) = 0$ for every $\sigma \in S_r$.
\end{proof}

\begin{proof}[Proof of the theorem]
Combine Proposition~\ref{prop:main} with Lemma~\ref{lem:reduce} and
the equivalence $[d]_t \mid f(t) \iff f(\zeta_d) = 0$ (valid since $d$
prime, $[d]_t = \Phi_d(t)$).
\end{proof}

\section{Range covered and remarks}

The theorem covers the following empirically verified cases of the
conjecture in \cite{PROVE-2026-07-31}:
\[
\begin{array}{lll}
r = 2, d = 3 & k = 2 & \text{(one case)}\\
r = 3, d = 5 & k = 2, 3, 4 & \text{(three cases)}\\
r = 4, d = 7 & k = 2, 3, 4, 5, 6 & \text{(five cases; probe verified $k=2,3$)}\\
r = 6, d = 11 & k = 2, \ldots, 10 & \text{(nine cases; probe verified $k=2$)}\\
r = 7, d = 13 & k = 2, \ldots, 12 & \text{(eleven cases; not probed)}\\
\vdots
\end{array}
\]

Note that for $r = 5$, $d = 9 = 3^2$ is \emph{not} prime; the theorem
does not directly apply, though the observation $m_d(\tilde\sigma h) = 0$
still holds (giving divisibility by $\Phi_9$; but $[9]_t = \Phi_3 \Phi_9$
also requires divisibility by $\Phi_3$, which the wreath cycle bound
does not immediately give).

\begin{remark}[Explicit verification]
The wreath-cycle claim of Corollary~\ref{cor:no-d-cycle} is
brute-force verified for $(r, k) \in \{(2,2), (3,2), (3,3), (3,4),
(4,2), (4,3)\}$ by enumerating all elements of $S_r \times S_k^r$ and
computing cycle types (script:
\texttt{\~{}/projects/probes/2026-07-31-r2-proof-verification/verify\_general\_r.py}).
Total wreath-element counts checked:
$8 + 48 + 1296 + 82944 + 384 + 31104 = 115784$; zero cycles of length
divisible by $d$ in every case.
\end{remark}

\begin{remark}[What remains]
The full conjecture asserts $[d]_t \mid m_\pi(t)$ iff
$k \bmod d \in \{2, \ldots, d-1\}$ for arbitrary $k$. Our theorem is
restricted to $k$ in the first ``allowed period,'' i.e., $k \in \{2,
\ldots, d-1\}$. The extension to $k \ge d$ (still with $k \bmod d \in
\{2, \ldots, d-1\}$) requires handling wreath elements with
$m_d(\tilde\sigma h) > 0$, where the individual Molien terms do
\emph{not} vanish at $\zeta_d$ and cancellation across $h$ becomes
essential. A natural attack is via the character of $\tilde\sigma h$ on
each fake-degree isotypic $S^\lambda$, using the Murnaghan--Nakayama
rule and the Springer regular-element criterion for
$(S_k \wr S_r, d)$-regular elements (compare Barcelo--Reiner 2005 for
wreath Springer theory).
\end{remark}

\begin{remark}[Composite $d$]
The prime hypothesis on $d = 2r - 1$ is convenient (allows the
$[d]_t \mid f \Leftrightarrow f(\zeta_d) = 0$ collapse); for composite
$d$, one must consider divisibility by each cyclotomic factor
$\Phi_e$, $e \mid d$, separately. The wreath-cycle bound
Corollary~\ref{cor:no-d-cycle} extends: if the argument at $\zeta_d$ uses
only ``$d$ prime, $r < d$,'' then for composite $d$ one replaces $d$ by
each prime $e \mid d$ (only useful when $e > r$).
\end{remark}

\bibliographystyle{alpha}
\begin{thebibliography}{9}
\bibitem{Wreath-general-r} C.\ Muse, \emph{The $r$-wreath extension of
Theorem D and a cyclotomic Hilbert-series identity}, note (2026-07-30).
\bibitem{PROVE-2026-07-31} \texttt{\~{}/state/PROVE.md}, period-$(2r-1)$
divisibility conjecture (2026-07-31).
\bibitem{Macdonald} I.G.~Macdonald, \emph{Symmetric Functions and Hall
Polynomials}, 2nd ed., Oxford Univ.~Press, 1995.
\bibitem{Barcelo-Reiner} H.~Barcelo, V.~Reiner, \emph{Springer's regular
elements over arbitrary fields}, Math.~Proc.~Cambridge Philos.~Soc. 141
(2006), 269--275.
\end{thebibliography}

\end{document}
